% Part: first-order-logic
% Chapter: sequent-calculus
% Section: proving-things-quant

\documentclass[../../../include/open-logic-section]{subfiles}

\begin{document}

\olfileid{fol}{seq}{prq}

\olsection{\usetoken{P}{derivation} با سورها}

\begin{ex}
یک $\Log{LK}$-!!{derivation} برای دنبالهٔ
$\lexists[x][\lnot !A(x)] \Sequent \lnot \lforall[x][!A(x)]$ به دست دهید.

هنگام کار با سورها باید مراقب باشیم شرط متغیر ویژه را نقض نکنیم، و
گاهی این امر ایجاب می‌کند ترتیب انجام برخی استنتاج‌ها را پس‌وپیش کنیم.
به‌طور کلی، بهتر است ابتدا به قواعد مشمول شرط متغیر ویژه بپردازیم
(این قواعد در برهان نهایی پایین‌تر قرار خواهند گرفت). همچنین بهتر است
کمی جلوتر را ببینیم و حدس بزنیم دنبالهٔ آغازین چه صورتی خواهد داشت.
در مورد ما، باید چیزی مانند $!A(a) \Sequent !A(a)$ باشد. یعنی هنگام
«وارونه‌کردن» قواعد سورها باید همان ترم---که آن را $a$ می‌نامیم---را
برای هر دو قاعدهٔ $\lforall$ و $\lexists$ برگزینیم. اگر برای
هر قاعده ترم متفاوتی برمی‌گزیدیم، به چیزی مانند
$!A(a) \Sequent !A(b)$ می‌رسیدیم که البته اشتقاق‌پذیر نیست.

طبق معمول، با نوشتن عبارت زیر آغاز می‌کنیم:
\begin{prooftree}
\AxiomC{}
\UnaryInf$\lexists[x][\lnot !A(x)] \fCenter \lnot \lforall[x][!A(x)]$
\end{prooftree}
می‌توانیم یا قاعدهٔ \LeftR{\exists} را اجرا کنیم یا قاعدهٔ
\RightR{\lnot} را. چون قاعدهٔ \LeftR{\exists} مشمول شرط متغیر ویژه
است، بهتر است زودتر به آن بپردازیم؛ پس آن را نخست اجرا می‌کنیم.
\begin{prooftree}
\AxiomC{}
\UnaryInf$ \lnot !A(a) \fCenter \lnot \lforall[x][!A(x)]$
\RightLabel{\LeftR{\lexists}}
\UnaryInf$ \lexists[x][\lnot !A(x)] \fCenter \lnot \lforall[x][!A(x)]$
\end{prooftree}
با اعمال وارونهٔ قواعد \LeftR{\lnot} و \RightR{\lnot} به عبارت زیر
می‌رسیم:
\begin{prooftree}
\AxiomC{}
\UnaryInf$\lforall[x][!A(x)] \fCenter !A(a)$
\RightLabel{\LeftR{\lnot}}
\UnaryInf$\lnot !A(a), \lforall[x][!A(x)] \fCenter $
\RightLabel{\LeftR{\Exchange}}
\UnaryInf$\lforall[x][!A(x)], \lnot !A(a) \fCenter $
\RightLabel{\RightR{\lnot}}
\UnaryInf$ \lnot !A(a) \fCenter \lnot \lforall[x] !A(x)$
\RightLabel{\LeftR{\lexists}}
\UnaryInf$ \lexists[x] \lnot !A(x) \fCenter \lnot \lforall[x] !A(x)$
\end{prooftree}
در این مرحله تنها گزینهٔ ما اجرای قاعدهٔ \LeftR{\forall} است. چون این
قاعده مشمول محدودیت متغیر ویژه نیست، مشکلی نداریم. به یاد آورید که
می‌خواهیم دنباله‌ای آغازین (به صورت $!A(a) \Sequent !A(a)$) به دست
آوریم؛ پس هنگام اعمال قاعده باید $a$ را به‌عنوان آرگومان $!A$ برگزینیم.
\begin{prooftree}
\Axiom$!A(a) \fCenter !A(a)$
\RightLabel{\LeftR{\lforall}}
\UnaryInf$\lforall[x][!A(x)] \fCenter !A(a)$
\RightLabel{\LeftR{\lnot}}
\UnaryInf$\lnot !A(a), \lforall[x][!A(x)] \fCenter $
\RightLabel{\LeftR{\Exchange}}
\UnaryInf$\lforall[x][!A(x)], \lnot !A(a) \fCenter $
\RightLabel{\RightR{\lnot}}
\UnaryInf$ \lnot !A(a) \fCenter \lnot \lforall[x][!A(x)]$
\RightLabel{\LeftR{\lexists}}
\UnaryInf$ \lexists[x][ \lnot !A(x)] \fCenter \lnot \lforall[x][!A(x)]$
\end{prooftree}
به‌ویژه هنگام کار با سورها مهم است که در این مرحله دوباره بررسی کنیم
شرط متغیر ویژه نقض نشده باشد. چون تنها قاعدهٔ اعمال‌شده‌ای که مشمول
شرط متغیر ویژه بود \LeftR{\exists} است و متغیر ویژهٔ~$a$ در دنبالهٔ
پایینیِ آن (دنبالهٔ پایانی) رخ نمی‌دهد، این !!{derivation} درست است.
\end{ex}

\begin{prob}
برای دنباله‌های زیر !!{derivation}sی لازم را به دست دهید:
\begin{enumerate}
\item $\Sequent (\lforall[x][!A(x)] \land \lforall[y][!B(y)]) \lif
\lforall[z][(!A(z) \land !B(z))]$.
\item $\Sequent (\lexists[x][!A(x)] \lor \lexists[y][!B(y)]) \lif
\lexists[z][(!A(z) \lor !B(z))]$.
\item $\lforall[x][(!A(x) \lif !B)] \Sequent \lexists[y][!A(y)] \lif !B$.
\item $\lforall[x][\lnot !A(x)] \Sequent \lnot\lexists[x][!A(x)]$.
\item $\Sequent \lnot\lexists[x][!A(x)] \lif \lforall[x][\lnot !A(x)]$.
\item $\Sequent \lnot\lexists[x][\lforall[y][((!A(x,y) \lif \lnot
!A(y,y)) \land (\lnot !A(y,y) \lif !A(x,y)))]]$.
\end{enumerate}
\end{prob}

\begin{prob}
برای دنباله‌های زیر !!{derivation}sی لازم را به دست دهید:
\begin{enumerate}
\item $\Sequent \lnot\lforall[x][!A(x)] \lif \lexists[x][\lnot!A(x)]$.
\item $(\lforall[x][!A(x)] \lif !B) \Sequent \lexists[y][(!A(y) \lif !B)]$.
\item $\Sequent \lexists[x][(!A(x) \lif \lforall[y][!A(y)])]$.
\end{enumerate}
(همهٔ این موارد به قاعدهٔ~$\RightR{\Contraction}$ نیاز دارند.)
\end{prob}

\end{document}
